% Generated draft: requires mathematical and source review.
\documentclass[11pt]{article}
\usepackage[margin=25mm]{geometry}
\usepackage{fontspec}
\setmainfont{Latin Modern Roman}
\usepackage{amsmath,amssymb,mathtools,unicode-math}
\setmathfont{Latin Modern Math}
\usepackage{enumitem,needspace,xurl,hyperref}
\setlength{\parindent}{0pt}
\setlength{\parskip}{5pt}
\setlength{\emergencystretch}{3em}
\newcommand{\N}{\mathbb{N}}
\newcommand{\Z}{\mathbb{Z}}
\newcommand{\Q}{\mathbb{Q}}
\newcommand{\R}{\mathbb{R}}
\newcommand{\C}{\mathbb{C}}
\newcommand{\F}{\mathbb{F}}
\newcommand{\A}{\mathbb{A}}
\newcommand{\PP}{\mathbb{P}}
\newcommand{\E}{\mathbb{E}}
\newcommand{\eps}{\varepsilon}
\newcommand{\dd}{\,\mathrm{d}}
\begin{document}
\section*{UnsolvedMath ID: 11000106}
\textbf{Source problem code:} catalogue entry
\subsection{Definitions and mathematical statement}
Let \(S,R\) be the closed surfaces of the source, with \(g(R)<g(S)\), and let \(\Gamma\leq\operatorname{Mod}_S\) have finite index. The stated conjecture is
\[
 \forall\rho:\Gamma\longrightarrow\operatorname{Mod}_R,
 \qquad[\operatorname{Mod}_R:\rho(\Gamma)]=\infty.
\]
Thus no such homomorphism has finite-index image. This is weaker than the further hope that \(\rho(\Gamma)\) is always finite; the source deliberately chooses the more cautious statement.

The allowed genus range is not explicit. Under the ordinary mapping-class convention, a closed sphere has trivial mapping class group, so allowing \(R=S^2\) contradicts the printed finite-index-image exclusion. Intended low-genus restrictions must therefore be recovered, not silently added.

The source relates the stronger hope to the unresolved finite-index first-cohomology question: if the finitely generated group \(\Gamma\) has infinite abelianization, it surjects onto \(\Z\), which can give infinite-image homomorphisms into a target mapping class group possessing an infinite-order element. Such maps alone do not give finite-index image in a non-virtually-cyclic target. The coefficients in the printed \(H^1(\Gamma)\) shorthand are omitted.

It reports W. Harvey and M. Korkmaz [HaK] proving the stated conjecture for \(\Gamma=\operatorname{Mod}_S\), under the intended source hypotheses. It says their finite-order-element methods do not extend to general finite-index subgroups. This is kept as a limitation of that approach, not a proof that no other method could work.
\subsection{Short English statement}
Can a finite-index mapping-class subgroup map with finite-index image to a lower-genus mapping class group? The printed conjecture excludes this, while finite image for every map is only a stronger hope and low-genus exclusions are missing.
\subsection{Sources}
\begin{itemize}
\item \emph{Problems on Mapping Class Groups and Related Topics}, edited by Farb, unnumbered Conjecture, supplied PDF page 88. \url{https://www.math.uchicago.edu/~farb/papers/mcgbook.pdf}. Recorded access 2026-07-29; extraction pdf-font-structure.
\item W. Harvey and M. Korkmaz [HaK], source attribution for the full-domain mapping-class-group case; complete bibliography and exact genus hypotheses are absent. The source also refers to its unresolved Problem 7 on \(H^1(\Gamma)\), whose complete statement is not supplied here.
\item Local UnsolvedMath registry, \texttt{data/unsolvedmath/problems.json}, permanent ID \(11000106\), exact source code \texttt{catalogue entry}; supplied OPEN-TRIAGE summary says truncated.
\end{itemize}
\end{document}
